% Part: normal-modal-logic
% Chapter: tableaux
% Section: simple-S5

\documentclass[../../../include/open-logic-section]{subfiles}

\begin{document}

\olfileid{nml}{tab}{s5}

\olsection{\usetoken{P}{tableau} Prasaja kanggo \Log{S5}}

\Log{S5} sah lan jangkep relatif marang kelas model universal, yaiku
model kang saben jagade bisa diakses saka saben jagad liyane. Ing model
universal, relasi aksesibilitas ora perlu dibedakake maneh:
``ana jagad~$w$ kang nyukupi $\mSat{M}{!A}[w]$'' bener yen lan mung
yen ana jagad $w$ mangkono kang bisa diakses saka~$u$. Dadi ing
\Log{S5}, model bisa ditepakake mung minangka himpunan jagad lan
valuasi~$V$. Iki nuduhake yen aturan !!{tableau} uga bisa
diprasajakake. Ing kasus umum, prefiks iku urutan bilangan bulat
positif supaya kita bisa nglacak prefiks endi kang maringi jeneng
jagad kang bisa diakses saka jagad liyane: $\sigma.n$ maringi jeneng
jagad kang bisa diakses saka~$\sigma$. Nanging ing \Log{S5} saben
jagad bisa diakses saka saben jagad liyane, mula pelacakan mangkono
ora perlu. Minangka gantine, kita bisa nggunakake bilangan bulat
positif minangka prefiks. Aturan kang wis diprasajakake ana ing
\olref{tab:rules-S5}.

\begin{table}
  \begin{center}
    \def\arraystretch{3}\def\fCenter{}
    \begin{tabular}{|c|c|}
    \hline
    \iftag{prvBox}{
      \AxiomC{\sFmla{\True}{\Box !A}[n]}
      \RightLabel{\TRule{\True}{\Box}}
      \UnaryInfC{\sFmla{\True}{!A}[m]}
      \DisplayProof
      &
      \AxiomC{\sFmla{\False}{\Box !A}[n]}
      \RightLabel{\TRule{\False}{\Box}}
      \UnaryInfC{\sFmla{\False}{!A}[m]}
      \DisplayProof\\
      $m$ wis dienggo & $m$ anyar
      \\[1ex]
      \hline}{}
    \iftag{prvDiamond}{
      \AxiomC{\sFmla{\True}{\Diamond !A}[n]}
      \RightLabel{\TRule{\True}{\Diamond}}
      \UnaryInfC{\sFmla{\True}{!A}[m]}
      \DisplayProof
      &
      \AxiomC{\sFmla{\False}{\Diamond !A}[n]}
      \RightLabel{\TRule{\False}{\Diamond}}
      \UnaryInfC{\sFmla{\False}{!A}[m]}
      \DisplayProof\\
      $m$ anyar & $m$ wis dienggo
      \\[1ex]
      \hline}{}
    \end{tabular}
  \end{center}
  \caption{Aturan kang diprasajakake kanggo \Log{S5}.}
  \ollabel{tab:rules-S5}
\end{table}

\begin{ex}
  Kita nyuguhake tableau katutup kang wis diprasajakake kanggo nuduhake
  $\Log{S5} \Proves \Ax{5}$, yaiku $\Diamond!A \lif \Box\Diamond!A$.
  \begin{oltableau}
    [\pFmla{\False}{\Diamond\formula{A} \lif \Box\Diamond \formula{A}}{1},
      just = \TAss
      [\pFmla{\True}{\Diamond \formula{A}}{1}, just = {\TRule{\False}{\lif}[1]}
        [\pFmla{\False}{\Box\Diamond \formula{A}}{1},
          just = {\TRule{\False}{\lif}[1]}
          [\pFmla{\False}{\Diamond \formula{A}}{2},
            just = {\TRule{\False}{\Box}[3]}
            [\pFmla{\True}{\formula{A}}{3},
              just = {\TRule{\True}{\Diamond}[2]}
                [\pFmla{\False}{\formula{A}}{3},
                  just = {\TRule{\False}{\Diamond}[4]}, close]
            ]
          ]
        ]
      ]
    ]
  \end{oltableau}
\end{ex}

\end{document}
